% Auto-generated by scripts/latex_synthesis.py -- do not hand-edit; regenerate from the
% source synthesis/<slug>.md instead, since edits here will be overwritten on the next run.
%
% Thesis-chapter style: report class (numeric \cite by default, no natbib/biblatex needed),
% standard margins, no title page/TOC -- this is meant to be read as a single chapter, not a
% standalone thesis. See scripts/latex_synthesis.py's CITATION REPLACEMENT STRATEGY comment
% for how the "Author et al. (Year)" text was turned into \cite{} calls below.
\documentclass[12pt]{report}
\usepackage[margin=1in]{geometry}
\usepackage[utf8]{inputenc}
\usepackage[T1]{fontenc}

\begin{document}

\chapter{is recurrence necessary for good sequence modeling performance}

% Contradiction/agreement analysis generated by scripts/synthesize.py.

Agreement: All three papers agree recurrence not strictly required for strong sequence modeling. \cite{1706.03762} drop recurrence and convolution entirely, get better BLEU than RNN-based systems at lower training cost. \cite{1810.04805} build BERT on pure Transformer (no recurrence), beat prior feature-based/fine-tuning approaches including ELMo (which use LSTMs).

Disagreement: \cite{1409.0473} do not challenge recurrence itself — they keep RNN encoder-decoder core, just add attention to fix fixed-length-vector bottleneck. So no direct disagreement on necessity-of-recurrence claim; Bahdanau et al. simply predates that question, still recurrence-based, while Vaswani et al. and Devlin et al. actively argue against needing it. Position, not contradiction: recurrence-optional (Vaswani, Devlin) vs recurrence-retained-but-improved (Bahdanau).

\bibliographystyle{plain}
\bibliography{is-recurrence-necessary-for-good-sequence-modeling-performan-contradictions}

\end{document}
